\subsection{Impedanz des Kondensators}
\label{subsec:impedanz-des-kondensators}

Die Impedanz eines Kondensators lässt sich direkt aus der Grundgleichung des
Kondensators herleiten. Für den Strom durch einen Kondensator gilt:

\[
i_C(t)=C\frac{du_C(t)}{dt}.
\]

Diese Gleichung zeigt bereits den entscheidenden Zusammenhang:
Der Strom durch den Kondensator ist proportional zur zeitlichen Änderung der
Kondensatorspannung.

Für eine sinusförmige Wechselspannung verwendet man in der Wechselstromrechnung
eine komplexe Darstellung. Man schreibt die Kondensatorspannung als

\[
\underline{u}_C(t)=\hat U_C e^{i\omega t}.
\]

Der reale physikalische Spannungsverlauf ist der Realteil dieser komplexen
Größe. Der Vorteil dieser Schreibweise liegt darin, dass die Ableitung nach der
Zeit besonders einfach wird:

\[
\frac{d\underline{u}_C(t)}{dt}
=
\frac{d}{dt}\left(\hat U_C e^{i\omega t}\right).
\]

Da $\hat U_C$ konstant ist, folgt:

\[
\frac{d\underline{u}_C(t)}{dt}
=
\hat U_C i\omega e^{i\omega t}.
\]

Damit erhält man:

\[
\frac{d\underline{u}_C(t)}{dt}
=
i\omega \hat U_C e^{i\omega t}.
\]

Da

\[
\underline{u}_C(t)=\hat U_C e^{i\omega t}
\]

gilt, kann man schreiben:

\[
\frac{d\underline{u}_C(t)}{dt}
=
i\omega \underline{u}_C(t).
\]

Setzt man dies in die Kondensatorgleichung ein, so erhält man:

\[
\underline{i}_C(t)
=
C\frac{d\underline{u}_C(t)}{dt}.
\]

Also:

\[
\underline{i}_C(t)
=
C\,i\omega \underline{u}_C(t).
\]

Damit folgt:

\[
\underline{i}_C(t)
=
i\omega C\,\underline{u}_C(t).
\]

Die Impedanz ist allgemein definiert als das Verhältnis von komplexer Spannung
zu komplexem Strom:

\[
Z_C=\frac{\underline{u}_C(t)}{\underline{i}_C(t)}.
\]

Aus

\[
\underline{i}_C(t)
=
i\omega C\,\underline{u}_C(t)
\]

folgt daher:

\[
Z_C
=
\frac{\underline{u}_C(t)}
{i\omega C\,\underline{u}_C(t)}.
\]

Die Spannung $\underline{u}_C(t)$ kürzt sich heraus:

\[
Z_C
=
\frac{1}{i\omega C}.
\]

Damit lautet die Impedanz des Kondensators:

\[
\boxed{
	Z_C=\frac{1}{i\omega C}
}
\]

Da für die imaginäre Einheit gilt:

\[
\frac{1}{i}=-i,
\]

kann man die Impedanz auch schreiben als:

\[
Z_C
=
-i\frac{1}{\omega C}.
\]

Also:

\[
\boxed{
	Z_C=-\,i\,\frac{1}{\omega C}
}
\]

In der Elektrotechnik verwendet man häufig den Buchstaben $j$ statt $i$, weil
$i$ bereits für den elektrischen Strom verwendet wird. Dann schreibt man:

\[
\boxed{
	Z_C=\frac{1}{j\omega C}
	=
	-j\frac{1}{\omega C}
}
\]

\subsubsection{Betrag und Phase}
\label{subsubsec:betrag-und-phase-kondensatorimpedanz}

Die Impedanz des Kondensators ist rein imaginär. Ihr Betrag ist:

\[
\left|Z_C\right|
=
\left|-i\frac{1}{\omega C}\right|.
\]

Da der Betrag von $i$ gleich $1$ ist, folgt:

\[
\boxed{
	\left|Z_C\right|=\frac{1}{\omega C}
}
\]

Der Betrag der Kondensatorimpedanz wird also mit steigender Frequenz kleiner.
Für kleine Frequenzen wirkt der Kondensator wie ein großer Widerstand, für hohe
Frequenzen wie ein kleiner Widerstand.

Die Phase von

\[
Z_C=-i\frac{1}{\omega C}
\]

beträgt

\[
\boxed{
	\varphi_C=-90^\circ
}
\]

bzw.

\[
\boxed{
	\varphi_C=-\frac{\pi}{2}
}
\]

Das bedeutet: Die Spannung am Kondensator hinkt dem Strom um $90^\circ$ nach.
Umgekehrt formuliert:

\[
\boxed{
	\text{Der Strom eilt der Kondensatorspannung um }90^\circ\text{ voraus.}
}
\]

\begin{DidacticBox}[Warum taucht $i\omega$ auf?]
	\label{box:warum-i-omega-auftaucht}
	
	Der zentrale mathematische Schritt ist die Ableitung der komplexen
	Exponentialfunktion:
	
	\[
	\frac{d}{dt}e^{i\omega t}
	=
	i\omega e^{i\omega t}.
	\]
	
	Eine zeitliche Ableitung entspricht in der komplexen Wechselstromrechnung also
	einer Multiplikation mit $i\omega$. Genau deshalb treten bei Kondensatoren und
	Spulen komplexe Impedanzen auf.
	
	Für den Kondensator führt diese Ableitung zu
	
	\[
	Z_C=\frac{1}{i\omega C}.
	\]
	
	Der Faktor $\frac{1}{i}=-i$ beschreibt dabei die Phasenverschiebung von
	$-90^\circ$ zwischen Kondensatorspannung und Strom.
\end{DidacticBox}